\documentclass[12pt]{article}

\usepackage[T2A]{fontenc}
\usepackage[utf8]{inputenc}
\usepackage[russian]{babel}

\usepackage[a4paper,vmargin=15mm,lmargin=30mm,rmargin=10mm]{geometry}

\usepackage{graphicx}
\usepackage{epstopdf}
\epstopdfsetup{suffix=}

\usepackage{array}

\pagestyle{empty}

\begin{document}

\begin{center}

Министерство науки и высшего образования Российской Федерации

\smallskip

ФГАОУ ВО <<УрФУ имени первого Президента России Б.Н. Ельцина>>

\medskip

\large

\textbf{ОТЗЫВ}

руководителя выпускной квалификационной работы

студента Корабельникова Алексея Алексеевича группы МЕНМ-280901
 
\medskip

\normalsize

\raggedright

\noindent
\textbf{Тема магистерской диссертации:} МНОГОМЕРНЫЙ АЛГОРИТМ ОВЫПУКЛЕНИЯ РОЯ ТОЧЕК, НАХОДЯЩИХСЯ В НЕОБЩЕМ ПОЛОЖЕНИИ.
\end{center}

Корабельников Алексей Алексеевич при работе над магистерской диссертацией проявил себя
следующим образом:

\noindent
\textbf{1.  Степень творчества:} высокая. При решении задачи выработал собственный подход к разработке алгоритма.

\smallskip

\noindent
\textbf{2.  Степень самостоятельности:} высокая. При работе над алгоритмом самостоятельно собрал и изучил итературы по существующим процедурам построения выпуклой оболочки роя в двумерном пространстве и в пространствах высокой размерности. Спроектировал алгоритм, нацеленный на ситуацию необщего положения точек овыпукляемого роя в многомерном пространстве. Разработал его программную реализацию на языке C\#.

\smallskip

\noindent
\textbf{3.  Работоспособность, прилежание, ритмичность:} демонстрировал высокую работоспособность и хорошее прилежание, можно также отметить ритмичность выполнения работы.

\smallskip

\noindent
\textbf{4.  Уровень специальной подготовки студента:} имеет отличную математическую и
компьютерную подготовку.

\smallskip

\noindent
\textbf{5. Возможность использования результатов в реальном секторе:} результаты планируется использовать в программах построения множеств разрешимости в выпуклых линейных задачах управления и дифференциальных играх.

\smallskip

\smallskip

\noindent
\textbf{6. Оценка соответствия ВКР требованиям ФГОС:} 

\centerline{
\begin{tabular}{|p{0.45\textwidth}|>{\centering}p{0.15\textwidth}|>{\centering}p{0.15\textwidth}|>{\centering}p{0.15\textwidth}|}
\hline
\centering Результаты освоения основной профессиональной образовательной программы, представленные в ВКР
&
Соответствует
&
В основном соответствует
&
Не соответствует
\cr
\hline
В части общекультурных компетенций & соответствует & & \cr
\hline
В части общепрофессиональных компетенций & соответствует & & \cr
\hline
В части профессиональных компетенций & соответствует & & \cr
\hline
\end{tabular}
}

\bigskip

\centerline{ОБЩЕЕ ЗАКЛЮЧЕНИЕ}

\medskip

В своей работе Корабельников Алексей Алексеевич разработал алгоритм построения выпуклой оболочки роя точек в многомерном пространстве, находящегося в необщем положении. Разработанный алгоритм был реализован в виде компьюбтерной программы на языке~C\#.

Магистерская диссертация Корабельникова Алексея Алексеевича полностью соответствует
требованиям, предъявляемым к выпускным квалификационным работам по направлению
09.04.03 Прикладная информатика, и заслуживает оценки <<отлично>>.

\medskip

Ф.И.О. руководителя ВКР: Кумков Сергей Сергеевич

\smallskip

Должность: доцент ДММиКН ИЕНиМ УрФУ

\smallskip

Уч. степень: кандидат физ.-мат. наук

\bigskip

\noindent
\textbf{Подпись} %
\raisebox{-1mm}{\hbox to 0pt{\rule{28mm}{0.4pt}}}% 
\hspace*{5mm}%
\raisebox{0mm}{\includegraphics[scale=0.75]{Kumkov.eps}}%
\hspace*{3mm}%
{\large /}% 
\raisebox{-1mm}{\hbox to 0pt{\rule{36mm}{0.4pt}}}% 
~~~~С.С.\,Кумков%
\hspace*{3cm}%
\textbf{Дата} 
\raisebox{-1mm}{\hbox to 0pt{\rule{4cm}{0.4pt}}}% 
~~<<\,03\,>> июня 2020 г.




 
\end{document}
